\section{Pilot Materials and Detailed Results}
\label{app:pilot-results}

This appendix presents the illustrative inputs, condition-level results and
representative outputs that support the pilot summary in
Section~\ref{sec:pilot}.

\subsection{Pilot Training Configuration}
\label{app:pilot-training}

The pilot parser used \texttt{google/gemma-2-9b-it} as its base model,
with LoRA fine-tuning on the fluent sentence--SBN pairs in PMB 5.1.0.
The training and development files were \texttt{train\_gold.json} and
\texttt{dev\_standard.json}, respectively. Tables~\ref{tab:pilot-lora}
and~\ref{tab:pilot-optimisation} record the training configuration used.

\begin{table}[htbp]
\centering
\small
\begin{tabular}{@{}p{0.36\textwidth}p{0.58\textwidth}@{}}
\toprule
Setting & Value \\
\midrule
Base model & \texttt{google/gemma-2-9b-it} \\
Precision & \texttt{bf16} \\
Attention implementation & \texttt{eager} \\
LoRA rank & 32 \\
LoRA scaling parameter ($\alpha$) & 64 \\
LoRA dropout & 0.05 \\
Bias adaptation & None (\texttt{bias = "none"}) \\
Target modules & \texttt{q\_proj}, \texttt{k\_proj}, \texttt{v\_proj},
  \texttt{o\_proj}, \texttt{gate\_proj}, \texttt{up\_proj},
  \texttt{down\_proj} \\
\bottomrule
\end{tabular}
\caption{Base-model and LoRA configuration. Adaptation covers both the listed
attention and MLP projections.}
\label{tab:pilot-lora}
\end{table}

\begin{table}[htbp]
\centering
\small
\begin{tabular}{@{}p{0.43\textwidth}p{0.51\textwidth}@{}}
\toprule
Setting & Value \\
\midrule
Learning rate & $2\times10^{-4}$ \\
Per-device batch size & 4 \\
Gradient accumulation steps & 4 \\
Reported effective batch size & 16 \\
Training epochs & 3.0 \\
Maximum sequence length & 1024 (updated from 512) \\
Warm-up ratio & 0.03 \\
Weight decay & 0.0 \\
Learning-rate scheduler & Cosine \\
Optimiser & \texttt{adamw\_torch} \\
Maximum gradient norm & 1.0 \\
Random seed & 42 \\
Gradient checkpointing & Enabled; \texttt{use\_reentrant = False} \\
Loss masking & Assistant SBN output only; prompt tokens excluded \\
Sequence packing & Disabled \\
Evaluation and checkpoint saving & Every 200 steps \\
Checkpoint retention limit & 3 \\
Final checkpoint selection & Load the best checkpoint by development loss
  (\texttt{eval\_loss}) \\
\bottomrule
\end{tabular}
\caption{Reported optimisation and training settings.}
\label{tab:pilot-optimisation}
\end{table}

\subsection{Input Conditions}

Table~\ref{tab:pilot-input-conditions} illustrates the fluent source sentence, three RSR
variants and their corresponding versions with an editing expression. Bold text
marks the inserted reparandum, underlined text marks the repair retained from
the source sentence, and italics mark the editing expression.

\begin{table}[ht]
\centering
\small
\begin{tabular}{@{}l>{\raggedright\arraybackslash}p{0.70\textwidth}@{}}
\toprule
\textbf{Input version} & \textbf{Example input} \\
\midrule
Fluent
  & She chatted with her friends about the baseball game. \\
RSR variant 1
  & \textbf{He}, \underline{she} chatted with her friends about the baseball
    game. \\
RSR variant 2
  & She chatted with her \textbf{coworkers}, \underline{friends} about the
    baseball game. \\
RSR variant 3
  & She chatted with her friends about the \textbf{basketball},
    \underline{baseball} game. \\
Variant 1 + editing expression
  & \textbf{He}, \textit{I mean}, \underline{she} chatted with her friends
    about the baseball game. \\
Variant 2 + editing expression
  & She chatted with her \textbf{coworkers}, \textit{I mean},
    \underline{friends} about the baseball game. \\
Variant 3 + editing expression
  & She chatted with her friends about the \textbf{basketball},
    \textit{I mean}, \underline{baseball} game. \\
\bottomrule
\end{tabular}
\caption{The seven input versions derived from one source sentence. The three
RSR variants are pooled within one condition in the main analysis; their
versions with an editing expression form the second RSR condition.}
\label{tab:pilot-input-conditions}
\end{table}

\subsection{Outcomes by Condition}

Table~\ref{tab:pilot-full} separates the two ways an output can fail to be
scored. An \emph{ill-formed} output cannot be read as SBN at all, whereas a
\emph{parse error} is produced when the sequence reads but yields an
unscoreable graph, typically through a cycle introduced by the model's own
index arithmetic.

\begin{table}[ht]
\centering
\small
\begin{tabular}{lrrrrrr}
\toprule
condition & $N$ & success \% & F1 (succ.) & F1 (failure-pen.) & ill-formed & parse error \\
\midrule
Fluent                     & 833 & 97.5 & 0.975 & 0.950 & 14 &   7 \\
\addlinespace
unmarked RSR, head         & 833 & 95.8 & 0.842 & 0.807 & 23 &  12 \\
unmarked RSR, mid          & 833 & 92.3 & 0.827 & 0.763 & 36 &  28 \\
unmarked RSR, tail         & 833 & 95.6 & 0.825 & 0.788 & 15 &  22 \\
\addlinespace
RSR + editing expression, head & 833 & 86.9 & 0.805 & 0.700 & 20 &  89 \\
RSR + editing expression, mid  & 833 & 71.8 & 0.779 & 0.559 & 30 & 205 \\
RSR + editing expression, tail & 833 & 70.7 & 0.803 & 0.568 & 14 & 230 \\
\bottomrule
\end{tabular}
\caption{Per-condition pilot results. F1 (succ.) averages over scoreable
outputs only; F1 (failure-pen.) assigns a failed output 0 and averages over all
inputs. The two rightmost columns count sentences.}
\label{tab:pilot-full}
\end{table}

The F1-among-successes and failure-penalised columns are both retained so that generation failure remains visible.

\subsection{Position Effects}

The unmarked-RSR and RSR-with-editing-expression groups are tested separately. Every sentence
appears at all three positions within its group, so the tests are paired:
Cochran's $Q$ with pairwise McNemar tests for parse success, and Friedman with
pairwise Wilcoxon tests for F1 among successes. Holm correction is applied to
the three pairwise comparisons within each group, at $\alpha = 0.05$. The F1
tests use complete triplets: 722 of 833 sentences for unmarked RSR and 408 for
the editing-expression conditions.

\begin{table}[ht]
\centering
\small
\begin{tabular}{llrlrr}
\toprule
group & outcome & $N$ & omnibus test & statistic & $p$ \\
\midrule
unmarked RSR         & parse success       & 833 & Cochran's $Q$ & 14.85 & 0.0006 \\
unmarked RSR         & F1 among successes  & 722 & Friedman      & 26.14 & $2.1\times10^{-6}$ \\
RSR + editing expression & parse success       & 833 & Cochran's $Q$ & 84.61 & $<10^{-6}$ \\
RSR + editing expression & F1 among successes  & 408 & Friedman      & 18.02 & 0.0001 \\
\bottomrule
\end{tabular}
\caption{Omnibus tests for an effect of reparandum position. All four are
significant at $\alpha = 0.05$.}
\label{tab:position-omnibus}
\end{table}

Table~\ref{tab:position-pairwise} gives the corresponding corrected pairwise
comparisons.

\begin{table}[ht]
\centering
\small
\begin{tabular}{llrrrl}
\toprule
 & & \multicolumn{3}{c}{Holm-corrected $p$} & \\
\cmidrule(lr){3-5}
group & outcome & head--mid & head--tail & mid--tail & separates \\
\midrule
unmarked RSR         & parse success      & 0.0037 & 0.897 & 0.0061 & mid (lower) \\
unmarked RSR         & F1                 & $1.9\times10^{-5}$ & $3.0\times10^{-6}$ & 0.439 & head (higher) \\
RSR + editing expression & parse success      & ${\sim}0$ & ${\sim}0$ & 0.633 & head (higher) \\
RSR + editing expression & F1                 & 0.0001 & 0.763 & 0.0005 & mid (lower) \\
\bottomrule
\end{tabular}
\caption{Pairwise position comparisons. McNemar for parse success, Wilcoxon
signed-rank for F1, Holm-corrected within each group.}
\label{tab:position-pairwise}
\end{table}

The F1 comparisons use only sentences with a scoreable output at all three
positions; this exclusion is reflected in the smaller $N$ values in
Table~\ref{tab:position-omnibus}.

\subsection{Representative Outputs}

\subsubsection{Editing-expression reification in the goat--cow example}
\label{app:pilot-goat-cow}

The example discussed in Section~\ref{sec:pilot} receives the following
output for the RSR input \textit{``A goat, I mean, a cow gave birth to a calf
with two heads''}:

\begin{lstlisting}[style=sbn]
% SPURIOUS: editing-expression material reified as content
goat.n.01                                    % goat
person.n.01 EQU speaker                      % I
mean.v.01 Theme -2 Agent -1 Time +1          % mean
time.n.08 EQU now
% GOLD-EQUIVALENT: the actual proposition
cow.n.01                                     % cow
give_birth.v.01 Agent -1 Time +1 Patient +2  % gave birth
time.n.08 TPR now                            % gave
calf.n.01 Part +2                            % calf
quantity.n.01 EQU 2                          % two
head.n.01 Quantity -1                        % heads
\end{lstlisting}

For comparison, the parser produces the following gold-equivalent output for
the fluent input \textit{``A cow gave birth to a calf with two heads''}:

\begin{lstlisting}[style=sbn]
cow.n.01                                     % cow
give_birth.v.01 Agent -1 Time +1 Patient +2  % gave birth
time.n.08 TPR now
calf.n.01 Part +2                            % calf
quantity.n.01 EQU 2                          % two
head.n.01 Quantity -1                        % heads
\end{lstlisting}

\subsubsection{Retention and loss of replacement candidates}

The following outputs document two error patterns behind the aggregate scores.
They are drawn from a purposive sample and are not used to estimate error
frequency. Both come from \textit{p07/d1400}, whose gold sentence is
\textit{``Is it true that you are in love with me?''} and whose fluent-input parse the
model recovers exactly:

\begin{lstlisting}[style=sbn]
time.n.08 EQU now
entity.n.01
true.a.01 Time -2 Proposition >1 EQU -1
    CONTINUATION <0
person.n.01 EQU hearer                        % you
time.n.08 EQU now
in_love.a.01 Experiencer -2 Time -1 Stimulus +1
person.n.01 EQU speaker                       % me
\end{lstlisting}

For the uncued input \textit{``Is it true that I am, you are in love with
me?''}, both candidates remain in the output:

\begin{lstlisting}[style=sbn]
% ... as above, up to the CONTINUATION
person.n.01 EQU speaker
time.n.08 EQU now
in_love.a.01 Experiencer -2 Time -1 Stimulus +1
person.n.01 EQU hearer
    CONTINUATION <1                           % the reparandum reading, retained
person.n.01 EQU hearer
time.n.08 EQU now
in_love.a.01 Experiencer -2 Time -1 Stimulus +1
person.n.01 EQU speaker
\end{lstlisting}

The output scores 0.743 against the fluent reference graph. When
\textit{``I mean''} is inserted, the output retains only the reparandum reading
and scores 0.926:

\begin{lstlisting}[style=sbn]
time.n.08 EQU now
entity.n.01
true.a.01 Time -2 Proposition >1 EQU -1
    CONTINUATION <0
person.n.01 EQU speaker                       % "I", the withdrawn candidate
time.n.08 EQU now
in_love.a.01 Experiencer -2 Time -1 Stimulus +1
person.n.01 EQU hearer                        % "me", now resolved to the hearer
\end{lstlisting}

These examples supplement the shorter output shown in
Section~\ref{sec:pilot}. Together, they show that high graph overlap can
coexist with retention of asserted reparandum material or loss of the accepted
repair.

\subsection{Recognition Probe Results}

That the parser fails to represent the RSR does not show the RSR was
invisible to it. A separate probe put the question to the base model, without
the adapter, which was asked to mark the reparandum, editing expression and repair in
plain text. Span scores are computed over sentences that were detected and
carry a true RSR.

\begin{table}[ht]
\centering
\small
\begin{tabular}{lrrrrrr}
\toprule
condition & $N$ & detected \% & off-template \% & scored & reparandum F1 & repair F1 \\
\midrule
unmarked RSR, head          & 833 & 30.7 &  0.0 & 254 & 0.667 & 0.487 \\
unmarked RSR, mid           & 833 & 39.0 &  0.0 & 324 & 0.503 & 0.420 \\
unmarked RSR, tail          & 833 & 64.2 &  0.7 & 536 & 0.489 & 0.340 \\
\addlinespace
RSR + editing expression, head & 833 & 100.0 &  0.6 & 831 & 0.900 & 0.862 \\
RSR + editing expression, mid  & 833 & 100.0 &  4.9 & 832 & 0.822 & 0.775 \\
RSR + editing expression, tail & 833 & 100.0 & 19.3 & 831 & 0.627 & 0.662 \\
\bottomrule
\end{tabular}
\caption{Base-model detection of self-repair in plain text. Off-template counts
outputs that did not follow the requested response format.}
\label{tab:ablation}
\end{table}

These figures report the condition-level values summarised in
Section~\ref{sec:pilot}. Detection rate and span F1 are reported separately
because identifying the presence of an RSR does not guarantee correct
segmentation of its components.

\section{Corpus Selection and Audit}
\label{app:corpus-selection}

This appendix records how structurally legal candidates are filtered and
sampled before model training. All boundaries and sampling decisions are
frozen before any main-experiment test output is inspected.

\subsection{Linguistic Calibration}

A 200-item calibration sample is drawn from the structurally legal pool,
stratified by noun, verb and adjective replacements and by the available
direct, anchor and movement constructions. Each item receives separate binary
judgements for semantic contrast and contextual plausibility. A replacement
such as \textit{``couch, I mean, sofa''} lacks sufficient contrast, whereas
\textit{``cake, I mean, carburettor, for his wedding''} lacks contextual
plausibility despite the lexical contrast.

Bidirectional NLI scores are calibrated against the contrast judgements, and
the mean token surprisal of the reparandum conditioned on its left context is
calibrated against the plausibility judgements. Each boundary, including the
upper surprisal threshold, is selected on the calibration sample and then held
fixed for the full candidate pool. Items that fail either criterion are
removed. After filtering, a separate 200-item sample receives the same manual
judgements. Its results estimate the retained corpus's rates of semantic
contrast, contextual plausibility and joint validity; this audit sample does
not alter the thresholds.

\subsection{Sampling and Split Audit}

All generated variants retain the PMB split and document identifier of their
source. A split audit verifies that no identifier occurs in more than one of
the training, development and test partitions. The final sample contains at
most two replacement sites from one source document. Where a document contains more
eligible sites, selection uses a fixed random seed and is stratified by part
of speech and construction. The realised corpus counts are frozen
before training and shared by the normalising and repair-aware conditions.

% 2026-10-24: preserve all base forms and append selected marked variants.
Every unmarked base form is retained. Marked variants are appended for
1,047 base items (30\%, rounded), including all 482 items flagged as
potentially ambiguous without a marker. The remaining 565 are sampled
proportionally to the non-ambiguous pool sizes by repair type, using largest
remainders and seed 20261024. Marker assignment uses seed 20261025 and
balances the four expressions within each type and across the selected set.
Ambiguity flags are heuristic; retaining the unmarked forms leaves textual
ambiguity as a limitation. Each surface variant has a unique identifier
and retains its base-item and source-document identifiers.
Test
sites retain their paired marked and unmarked forms under the same source
identifier, allowing them to be aggregated before statistical comparison.

\subsection{Challenge-Set Selection}

The non-RSR hold-out begins with 105 candidate PMB training documents,
which are removed from both the fluent training component and the RSR source
pool before fine-tuning. This leaves 9,447 PMB training documents. Manual
screening retains lists, appositions and coordination structures that resemble
the generated RSR surface pattern but carry no replacement reading. Comparable
items from the standard PMB test partition are added without changing their
fluent gold SBNs.

The discourse challenge set is assembled after the main sentence-level data
are frozen. Selected PMB documents and hand-written cases cover fluent
multi-sentence controls, later reference to accepted content, reference to
both replacement candidates and correction between discourse turns. Only
referents licensed as accessible by Section~\ref{sec:paired-configuration} are
used in targets that refer back to withdrawn content.

\section{Model and Evaluation Specifications}
\label{app:model-evaluation}

This appendix records the implementation and scoring decisions for the model
experiment in Section~\ref{sec:model}. The main text states the
comparisons and the evidence needed to answer the research question; the
details below make those decisions reproducible.

\subsection{Model Configurations and Run Matrix}
\label{app:model-config}

Each of the 2B, 9B and 27B models receives three matched-seed runs under each
condition: fluent only, RSR normalisation and repair-aware. The full design
therefore contains 27 runs. Matching seed values across conditions within each
model size reduces avoidable variation in the normalising--repair-aware
comparison.

Every run uses the same parsing instruction and LoRA adaptation method. The
prompt asks for an SBN sequence but does not label the input as fluent or
RSR-containing. Within each model size, the fluent data, RSR inputs, training steps
and optimisation settings are held constant wherever the conditions permit.
The development checkpoint with the lowest development loss is selected
without consulting test outputs. The final experiment record will include the
model checkpoint, prompt, LoRA targets, rank and scaling, optimiser settings,
batch construction, sequence limit, seed and selected checkpoint for every
run.

\subsection{Repair-Aware Evaluation Metrics}
\label{app:model-metrics}

RSR-structure accuracy is a binary score for each generated RSR item. A
prediction receives one point only when all four requirements hold:

\begin{enumerate}[nosep]
  \item the SBN sequence can be converted into a well-formed graph;
  \item the output contains the paired \sbnterm{CORRECTION} and
        \sbnterm{CONJUNCTION} contexts required by the target;
  \item the target reparandum is in the withdrawn context and the target repair
        is in the accepted context; and
  \item both parts have the semantic-role links to shared content required by
        the gold graph.
\end{enumerate}

The four requirements are also reported separately. These component scores
distinguish a missing configuration from reversed status or incorrect
attachment. Failure-penalised SMATCH++ assigns zero to an output that cannot be
scored and otherwise measures overlap with the complete repair-aware graph.
Results grouped by the direct, anchor and movement constructions show whether
the construction used affects parsing difficulty.

Accepted-meaning evaluation projects a prediction onto the speaker's final
meaning. The withdrawn context and the RSR-specific operators are removed, while
the accepted content and the surrounding graph are retained. The projected
graph is compared with the fluent gold PMB graph. On the unmodified PMB
test partition, the evaluation reports failure-penalised SMATCH++, graph
validity and the proportion of outputs that introduce an RSR configuration.

\subsection{Worked Scoring Example}
\label{app:model-scoring}

Consider \textit{``Let's stand, I mean, sit in the shade''}. The target marks
\textit{stand} as withdrawn and \textit{sit} as accepted, with both events
connected to the shared participants and time. An output that places both
verbs in ordinary asserted content may retain much of the target graph and
receive partial SMATCH++ credit. Its RSR-structure accuracy is nevertheless
zero because it does not express the replacement. An output receives full
RSR-structure credit only when the contexts, status assignments and shared
links are all correct. Removing its withdrawn context should then yield the
fluent meaning of \textit{``Let's sit in the shade''} for accepted-meaning
evaluation.

\subsection{Error Annotation Guidelines}
\label{app:model-errors}

The diagnostic categories are fixed before the main test outputs are
inspected. They distinguish a missing RSR configuration, reversed
withdrawn and accepted status, an incorrect link to shared content, a
reparandum left in accepted content, an editing expression represented as ordinary
semantic content, a spurious RSR analysis on fluent input and malformed SBN. Categories
may co-occur when they describe separate faults in one output. Any pattern
first identified from the main test outputs will be reported as a post-hoc
observation and excluded from the planned comparison.

\subsection{Statistical Comparisons}
\label{app:model-statistics}

All variants derived from one PMB source document are averaged before a
comparison, making the source document the unit of observation. For each model
size and matched seed, paired document-level differences are reported with
95\% confidence intervals from 10,000 paired bootstrap samples. Within each
model size, the mean and range of the three seed-specific differences
summarise training variation. A result is interpreted as stable only when all
three differences have the same direction. Holm correction is applied within
each family of secondary comparisons.

\section{A Later-Reference Challenge Example}
\label{app:later-reference}

The following hand-written target extends the Berlin--Dublin replacement from
the introduction with a second utterance that refers to both cities. It
illustrates the type of later-reference case reserved for the test-only
challenge set in Section~\ref{sec:data}.

\begin{lstlisting}[style=sbn]
% Utterance 1: What's the weather like in Berlin, I mean, in Dublin?
entity.n.01 EQU ?                              % what
be.v.01 Theme -1 Time +1 Co-Theme +2           % 's ... like
time.n.08 EQU now
weather.n.01                                   % the weather
    CORRECTION <1
city.n.01 Name "Berlin" LocationOf -1          % [reparandum] in Berlin
    CONJUNCTION <2                             % I mean,
city.n.01 Name "Dublin" LocationOf -2          % [repair] in Dublin
% Utterance 2: Show me the airports of these two cities.
    CONTINUATION <1
person.n.01 EQU hearer                         % imperative addressee
show.v.04 Agent -1 Recipient +1 Theme +2       % show
person.n.01 EQU speaker                        % me
airport.n.01 Location +2                       % the airports of
quantity.n.01 EQU 2                            % two
city.n.01 ANA -7 ANA -6 Quantity -1            % these ... cities
\end{lstlisting}

Here, \sbnterm{CONTINUATION <1} opens the context for the second
utterance and links it to the accepted context. The two
\sbnterm{ANA} links encode Berlin and Dublin as the intended plural
antecedent. This hand-written target operationalises the challenge case.

\section{RSR Constructions and Transformation Coverage}
\label{app:devices}

This appendix details the three constructions introduced in
Section~\ref{sec:paired-configuration}, states which one applies to a given
site, and reports the coverage of the concept-replacement operator over the PMB
training partition.

\subsection{The anchor construction}
\label{app:devices-anchor}

The paired configuration attaches both replacement candidates to the content
they share. Where the reparandum carries a role of its own, as
\sbnterm{Berlin} does in Figure~\ref{fig:weather-rsr-comparison}(b), that role
does the work. A subject concept selected as the repair has no such role. Taking \textit{``Josh
bought a cake''} as the source,

\begin{lstlisting}[style=sbn]
male.n.02 Name "Josh"                  % Josh
time.n.08 TPR now                      % (past)
buy.v.01 Agent -2 Time -1 Theme +1     % bought
cake.n.01                              % a cake
\end{lstlisting}

\noindent moving \sbnterm{male.n.02} into a \sbnterm{CORRECTION} context would
leave it with no link at all, since its only connection to the sequence is the
\sbnterm{Agent} role that the buying event points at it. The anchor construction inserts
a neutral \sbnterm{entity.n.01} concept before the \sbnterm{CORRECTION},
redirects the inbound role onto it, and attaches both candidates with
\sbnterm{EQU}:

\begin{lstlisting}[style=sbn]
entity.n.01                            % (anchor)
    CORRECTION <1
male.n.02 Name "Josh" EQU -1           % Josh
    CONJUNCTION <2
female.n.02 Name "Mary" EQU -2         % I mean, Mary
time.n.08 TPR now
buy.v.01 Agent -4 Time -1 Theme +1     % bought
cake.n.01                              % a cake
\end{lstlisting}

The buying event now takes the anchor as its \sbnterm{Agent}, and the two
candidates each assert an identity with it from inside their own context. The
withdrawn content is the identification of that individual as Josh, and the
accepted content is the identification as Mary. Neither \sbnterm{EQU} link
crosses a boundary in a forbidden direction, since both are negative indices
reaching outward to a concept in the parent context.

\subsection{The movement construction}
\label{app:devices-movement}

A positive index may not cross from a \sbnterm{CORRECTION} context into the
repair that follows it, which creates a problem whenever the repair concept
points forward. In \textit{``Tom broke the window''},

\begin{lstlisting}[style=sbn]
male.n.02 Name "Tom"                   % Tom
time.n.08 TPR now                      % (past)
break.v.01 Agent -2 Time -1 Theme +1   % broke
window.n.01                            % the window
\end{lstlisting}

\noindent a reparandum copy of \sbnterm{break.v.01} placed before the
\sbnterm{CORRECTION} would need \sbnterm{Theme +1} to reach
\sbnterm{window.n.01}, and that index would run through both separators. The
window is moved in front of the \sbnterm{CORRECTION} instead, so that both
copies reach it backwards:

\begin{lstlisting}[style=sbn]
male.n.02 Name "Tom"                   % Tom
time.n.08 TPR now
window.n.01                            % (moved from the end)
    CORRECTION <1
open.v.01 Agent -3 Time -2 Theme -1    % opened
    CONJUNCTION <2
break.v.01 Agent -4 Time -3 Theme -2   % I mean, broke
\end{lstlisting}

The movement construction changes the order of the representation:
\textit{the window} is uttered last but represented third. This is the
departure from word order noted in
Section~\ref{sec:formalism-coverage}, and it is accepted because the
alternative is to quarantine material the speaker kept. The move is confined to
arguments of the repair concept itself. Material that was uttered and
committed between the reparandum and the repair, as in the Lidl example, would
have to be carried past both separators, and that case is left outside the
configuration.

\subsection{Which construction applies}
\label{app:devices-order}

The direct construction is tried first at every site, since it changes neither
the inventory of concepts nor their order. Where it fails, the choice between
the anchor and movement constructions follows from the type of the repair
concept. The anchor construction asserts
\sbnterm{entity.n.01 EQU}~$c$ for each candidate $c$, which reads correctly
only when $c$ denotes a referent. Equating an entity with a property or an
event is not what the subject-RSR case supports, so for adjectives, adverbs
and verbs the movement construction is tried before the anchor, even though it
costs surface order. For nouns the order is reversed.
Table~\ref{tab:devices-coverage} shows the resulting distribution. Nouns take
the anchor construction at about the rate they take the direct construction,
while nearly all of the movement cases are adjectives.

\subsection{Coverage by part of speech}
\label{app:devices-coverage}

Table~\ref{tab:devices-coverage} breaks down the funnel of
Table~\ref{tab:coverage-funnel} by the part of speech of the repair concept
and by the construction the site ended up needing.

\begin{table}[ht]
\centering
\small
\begin{tabular}{@{}lrrrrrrr@{}}
\toprule
& & & & & \multicolumn{3}{c}{Construction used} \\
\cmidrule(l){6-8}
POS & Lexical & Aligned & Usable & Legal & Direct & Anchor & Movement \\
\midrule
Noun      & 11,781 & 10,060 & 9,560 & 9,525 & 4,576 & 4,949 &   0 \\
Verb      &  8,609 &  5,091 & 4,229 & 4,221 & 4,170 &     7 &  44 \\
Adjective &  2,842 &  2,735 & 2,692 & 2,679 & 2,050 &   148 & 481 \\
\midrule
Total     & 23,232 & 17,886 & 16,481 & 16,425 & 10,796 & 5,104 & 525 \\
\bottomrule
\end{tabular}
\caption{Candidate sites and transformation outcomes over the 9,552 English
gold documents of the PMB training partition. \emph{Aligned} counts concepts
that can be tied to an uttered word, \emph{usable} those for which WordNet
contains an eligible reparandum candidate, and \emph{legal} those where the transformation produces a
well-formed graph.}
\label{tab:devices-coverage}
\end{table}

Two of the drops deserve comment. Some 23\% of lexical concepts cannot be
aligned to an uttered word, and verbs account for most of that loss; a concept
with no surface expression cannot be given a competing expression either. A
further 6\% survive alignment but have no WordNet reparandum available, which
mainly affects senses with no co-hyponyms of the right part of speech.

\subsection{Why transformations fail}
\label{app:devices-blockers}

Of the 16,481 usable sites, 56 produce no legal graph under any of the three
options. Three conditions account for them, and they overlap, since a site may
be rejected differently by each option tried:

\begin{itemize}[nosep]
  \item \textbf{Index overflow} (45 sites). The PMB's released index pattern
    matches a single digit, so a role index may not exceed nine positions in
    either direction. Inserting two contexts pushes some existing indices past
    that ceiling. The limit belongs to the released tooling rather than to the
    notation, and it would disappear if the pattern were widened.
  \item \textbf{No attachment for the reparandum} (23 sites). Neither the
    original role nor the anchor construction leaves the withdrawn concept
    connected to the rest of the graph.
  \item \textbf{Movement unavailable} (50 sites). The repair concept has
    forward arguments that cannot be brought in front of the
    \sbnterm{CORRECTION} without moving material that belongs elsewhere.
\end{itemize}

Two further checks, for cycles and for collisions between the inserted
separators and those already present, reject individual attempts rather than
whole sites. Every site that triggered one was rescued by another option, so
neither appears in the tally above.

\subsection{Example Generated Corpus Record}
\label{app:generated-record}

Table~\ref{tab:worked-record} shows the complete record stored for one
generated replacement. The fields retain the PMB source, the choices made by
the transformation operator, both surface forms, the repair-aware SBN and the
outputs used for structural and linguistic validation. They are grouped by
function so that the generated training pair can be read separately from the
metadata retained for tracing and filtering it.

\begin{table}[ht]
\centering
\small
\begin{tabular}{@{}lp{0.62\textwidth}@{}}
\toprule
Field & Value \\
\midrule
\multicolumn{2}{@{}l}{\textit{PMB source}} \\
\addlinespace[0.25em]
\texttt{doc\_id}             & \texttt{p42/d1468} \\
\texttt{split}               & \texttt{train} \\
\texttt{nl\_clean}           & Let's sit in the shade. \\
\texttt{sbn\_clean}          & \texttt{time.n.08 TSU now person.n.01 EQU speaker sit.v.01 Time -2 Theme -1 Location +1 shade.n.01} \\
\midrule
\multicolumn{2}{@{}l}{\textit{Repair selection and transformation}} \\
\addlinespace[0.25em]
\texttt{concept\_pos}        & \texttt{2} (the repair concept is the third in
                               the sequence; positions are counted from zero) \\
\texttt{repair\_synset}      & \texttt{sit.v.01} \\
\texttt{reparandum\_synset}  & \texttt{stand.v.01} \\
\texttt{reparandum\_surface} & stand \\
\texttt{strategy}            & \texttt{inplace} \\
\texttt{devices}             & \texttt{d1\_drop\_role} \\
\texttt{blockers}            & --- \\
\texttt{max\_abs\_index}     & \texttt{3} (the largest role index in the output,
                               taken without its sign; the ceiling is nine) \\
\midrule
\multicolumn{2}{@{}l}{\textit{Generated pair}} \\
\addlinespace[0.25em]
\texttt{nl\_repair}          & Let's stand, sit in the shade. \\
\texttt{nl\_repair\_interr}  & Let's stand, I mean, sit in the shade. \\
\texttt{sbn\_repair}         & \texttt{time.n.08 TSU now person.n.01 EQU speaker CORRECTION <1 stand.v.01 Time -2 Theme -1 CONJUNCTION <2 sit.v.01 Time -3 Theme -2 Location +1 shade.n.01} \\
\midrule
\multicolumn{2}{@{}l}{\textit{Validation and filtering}} \\
\addlinespace[0.25em]
\texttt{legal}               & \texttt{1}, meaning the transformation produced a
                               well-formed graph. Items recorded as \texttt{0}
                               are dropped before any further filtering \\
\texttt{entail\_fwd}         & 0.0005 \\
\texttt{entail\_bwd}         & 0.0005 \\
\texttt{synonymy}            & 0.0005 \\
\bottomrule
\end{tabular}
\caption{A complete record produced by the current corpus-construction
pipeline. The ceiling on \texttt{max\_abs\_index} is the single-digit index
pattern described in Appendix~\ref{app:devices-blockers}; a transformation
that exceeds it is recorded as \texttt{legal 0} with
\texttt{index\_overflow} among its blockers.}
\label{tab:worked-record}
\end{table}